\documentclass{article}
\usepackage{propositions}

%% The `enum' style and the `enumprop' counter.
%%
%% enum is `numbered' with a counter that restarts, so that a list numbers from
%% one each time it is used, the way enumerate does.  numpropi, which
%% `numbered' uses at level one, runs on through the document instead -- right
%% for propositions referred to from elsewhere, wrong for a one-off list.
%%
%% The counter is reset at every outermost prop or inlineprop environment, so
%% nesting must not reset it: a sublist is not an outer environment.

\begin{document}

\section{A one-off list restarts; \texttt{numbered} does not}

\begin{prop}[style=enum]
  \item First. \label{e:a}
  \item Second. \label{e:b}
\end{prop}

\begin{prop}[style=enum]
  \item Must restart at one. \label{e:c}
  \item Second again. \label{e:d}
\end{prop}

References: \ref{e:a}, \ref{e:b}, \ref{e:c}, \ref{e:d}.  Expected:
(1), (2), (1), (2) --- each list numbered independently.

The default \texttt{numbered} style, for contrast, carries on:

\begin{prop}
  \item Runs on. \label{n:a}
\end{prop}

\begin{prop}
  \item Still running. \label{n:b}
\end{prop}

References: \ref{n:a}, \ref{n:b}.  Expected: two consecutive numbers, not
\texttt{(1)} twice.

\section{Nesting does not restart the counter}

A sublist is not an outermost environment, so the outer numbering must
continue across it, and the inner levels take the usual sub-item styles.

\begin{prop}[style=enum]
  \item First. \label{e:n1}
  \begin{prop}
    \item Sub-item. \label{e:n1a}
    \item Another sub-item. \label{e:n1b}
  \end{prop}
  \item Second, after a sublist. \label{e:n2}
  \item Third. \label{e:n3}
\end{prop}

The \emph{labels} must read (1), a., b., (2), (3): the outer count is
unaffected by the sublist, and the sub-items take the level-two style.

References: \ref{e:n1}, \ref{e:n1a}, \ref{e:n1b}, \ref{e:n2}, \ref{e:n3}.
Expected: (1), (1a), (1b), (2), (3) --- a sub-item's \emph{reference} carries
its parent's number, which is what \texttt{leveltwo} is for; only the label
is the bare letter.

\section{The restart survives an intervening list}

\begin{prop}[style=enum]
  \item One. \label{e:x1}
  \item Two. \label{e:x2}
\end{prop}

\begin{prop}
  \item An ordinary numbered list in between.
\end{prop}

\begin{prop}[style=enum]
  \item Must be one again. \label{e:y1}
\end{prop}

References: \ref{e:x1}, \ref{e:x2}, \ref{e:y1}.  Expected: (1), (2), (1).

\section{\texttt{inlineprop} counts as an outermost environment}

\begin{prop}[style=enum]
  \item One. \label{e:i1}
\end{prop}

\begin{inlineprop}[style=enum]
  \item An inline item. \label{e:i2}
\end{inlineprop}

\begin{prop}[style=enum]
  \item One again. \label{e:i3}
\end{prop}

References: \ref{e:i1}, \ref{e:i2}, \ref{e:i3}.  Expected: (1), (1), (1) ---
each is an outermost environment, so each restarts.

\section{Starting at a chosen number}

The reset happens at \verb|\begin{prop}|, so setting the counter just before
the environment has no effect --- the reset overwrites it.

\begingroup
\setcounter{enumprop}{40}
\begin{prop}[style=enum]
  \item Set before the environment. \label{e:s1}
\end{prop}
\endgroup

To start a list elsewhere than one, set the counter \emph{inside} the
environment, after the reset has run:

\begingroup
\begin{prop}[style=enum]
  \setcounter{enumprop}{40}
  \item Set inside the environment. \label{e:s2}
  \item And the next. \label{e:s3}
\end{prop}
\endgroup

References: \ref{e:s1}, \ref{e:s2}, \ref{e:s3}.  Expected: (1), (41), (42).

\end{document}
